\scanpage{19}
\uline{Théorème}
\margin{\unsure{Grotesque|Grothendieck|Ceci lorsque}, en fait on prouve déjà que $\uline{\mathrm{Pic}}^0_{X/S}$ est \uline{ouvert}, donc en l'occurrence \uline{simple} sur $S$.}
Soient $f:X\unsure{-|\to}S$ un morphisme
% slanted
\margin{Prouvé seulement en \uuline{car. $0$} pour l'instant\ldots.}
projectif, \deleted{plat, \inserted{et} à fibres géométriquement \deleted{\unsure{connexes}} \inserted{intègres} et
simples en codim $1$}, \inserted{\unsure{cf. page|cf. prop.} \unsure{2|?}} la base étant loc.
noethérienne et \uline{réduite}.
% slanted
\margin{est-ce nécessaire ? Non, grâce \unsure{à la simplicité|la rigidité} des Modules formels des V.A. $\ast$}
Alors il existe un
sous-préschéma abélien unique $A$ de $\uline{\mathrm{Pic}}_{X/S}$,
dont l'ensemble sous-jacent soit la réunion
des composantes connexes des fibres de $\uline{\mathrm{Pic}}_{X/S}$.

Soit en effet $P$ la dite réunion. Comme
ses fibres sont \uline{propres} sur $k$ [\uuline{\unsure{réf|vif|vrai}} \ldots],
il s'ensuit \deleted{\unsure{que}} par un théorème connu
[renvoyer aux critères de propreté au Ch. II]
que $P$ est \uline{propre} sur $S$ [question:
en est-il de $=$ de $\mathrm{Pic}^{\tau}_{X/S}$ ? De façon
générale, est-il vrai que $\uline{\mathrm{Pic}}_{X/S}$ est
réunion disjointe d'ouverts \uline{propres} sur
% slanted
\margin{? Non si on suppose $\unsure{X|S}$ seulement formel\ldots.}
$S$ ?]. Nous cherchons un sous-schéma
fermé de $P$, qui soit plat sur $S$ et
qui soit un sous-préschéma en groupes
de $\uline{\mathrm{Pic}}_{X/S}$ \inserted{et abélien sur $S$, \unsure{ensemblt égl à|ensemble ss-jt:|ensemblistement égal à} $\unsure{P|X}$}. \deleted{\unsure{C'est le}} Soit, pour $S'/S$
variable, $F(S')=$ \unsure{ens|les|l'ens.} des sous-schémas
fermés de $P_{S'}$, qui sont des sous-schémas
en groupes de $G$, \inserted{abéliens sur $S$} ensemblistiquement
égaux à $P_{S'}$. Je dis que \unsure{l'on|l'ens} $F(S')$
% archive sheet number: 18
